\chapter{2026年上海卷（春）}

\section{填空题}

\begin{problem}
已知集合 \(A=\{2,4\}\)，\(B=\{2,3,m\}\)，若 \(A\subseteq B\)，则 \(m=\fillinblank{}\).
\end{problem}

\begin{answer}
\(4\)
\end{answer}

\begin{solution}
由 \(A\subseteq B\) 知，集合 \(A\) 中的元素 \(4\) 也属于 \(B\)，故 \(m=4\).
\end{solution}

\begin{problem}
关于 \(x\) 的不等式 \(\dfrac{x+2}{x-3}<0\) 的解集为\fillinblank{}.
\end{problem}

\begin{answer}
\((-2,3)\)
\end{answer}

\begin{solution}
由 \(\dfrac{x+2}{x-3}<0\) 得 \((x+2)(x-3)<0\)，解得 \(-2<x<3\). 故解集为 \((-2,3)\).
\end{solution}

\begin{problem}
已知 \(\bs{a}=(x,3)\)，\(\bs{b}=(4,6)\)，若 \(\bs{a}\parallel\bs{b}\)，则 \(x=\fillinblank{}\).
\end{problem}

\begin{answer}
\(2\)
\end{answer}

\begin{solution}
由向量平行可得 \(6x-4\times3=0\)，即 \(6x-12=0\)，解得 \(x=2\).
\end{solution}

\begin{problem}
在平面直角坐标系中，点 \((1,2)\) 到直线 \(4x-3y+5=0\) 的距离为\fillinblank{}.
\end{problem}

\begin{answer}
\(\dfrac35\)
\end{answer}

\begin{solution}
由点到直线的距离公式，
\(d=\frac{|4\cdot1-3\cdot2+5|}{\sqrt{4^2+(-3)^2}}=\frac{3}{5}\).
因此点到直线的距离为 \(\dfrac35\).
\end{solution}

\begin{problem}
在二项式 \(\left(\dfrac1x+3x^2\right)^6\) 的展开式中，\(x^{-3}\) 的系数为\fillinblank{}.
\end{problem}

\begin{answer}
\(18\)
\end{answer}

\begin{solution}
通项为
\(T_{k+1}=\mathrm C_6^k\left(\frac1x\right)^{6-k}(3x^2)^k=\mathrm C_6^k3^kx^{3k-6}\).
令 \(3k-6=-3\)，得 \(k=1\). 故所求系数为 \(\mathrm C_6^1\cdot3=18\).
\end{solution}

\begin{problem}
若 \(a>0\)，\(b>0\)，且 \(a+2b=4\)，则 \(ab\) 的最大值为\fillinblank{}.
\end{problem}

\begin{answer}
\(2\)
\end{answer}

\begin{solution}
由 \(a+2b=4\) 得 \(a=4-2b\). 于是
\(ab=b(4-2b)=-2\left(b-1\right)^2+2\)，
故 \(ab\) 的最大值为 \(2\)，当 \(b=1\)，\(a=2\) 时取到.
\end{solution}

\begin{problem}
从\(5\text{ 人}\)中选\(3\text{ 人}\)去演讲，若甲一定去，则共有\fillinblank{}\(\text{ 种}\)选法.
\end{problem}

\begin{answer}
\(6\)
\end{answer}

\begin{solution}
甲已确定参加，只需从其余\(4\text{ 人}\)中再选\(2\text{ 人}\)，故共有 \(\mathrm C_4^2=6\text{ 种}\).
\end{solution}

\begin{problem}
已知点 \(P\) 在抛物线 \(y^2=4x\) 上，且点 \(P\) 到焦点的距离是到 \(y\) 轴距离的两倍，则点 \(P\) 的横坐标为\fillinblank{}.
\end{problem}

\begin{answer}
\(1\)
\end{answer}

\begin{solution}
设 \(P(x_0,y_0)\)（\(x_0\ge0\)）. 抛物线 \(y^2=4x\) 的焦点为 \((1,0)\)，准线为 \(x=-1\). 由题意，
\(x_0+1=2x_0\)，
故 \(x_0=1\).
\end{solution}

\begin{problem}
已知 \(m>1\)，对于所有满足 \(|z|=2\) 的复数 \(z\)，\(|z-\i|\) 的最小值与 \(|z-m|\) 的最小值相同，则 \(m=\fillinblank{}\).
\end{problem}

\begin{answer}
\(3\)
\end{answer}

\begin{solution}
\(|z|=2\) 表示复数 \(z\) 对应的点在以原点为圆心，半径为 \(2\) 的圆上. 点 \((0,1)\) 到圆上点的最小距离为 \(2-1=1\). 当 \(1<m<2\) 时，点 \((m,0)\) 到圆上点的最小距离为 \(2-m\)，由 \(2-m=1\) 得 \(m=1\)，与 \(m>1\) 矛盾；当 \(m\ge2\) 时，最小距离为 \(m-2\)，由 \(m-2=1\) 得 \(m=3\).
\end{solution}

\begin{problem}
在 \(\triangle ABC\) 中，\(D\)，\(E\) 在边 \(BC\) 上，且 \(BD=DE=EC\)，\(AD=1\)，\(\angle DAE=\dfrac{\pi}{3}\)，则 \(\overrightarrow{AB}\cdot\overrightarrow{AC}\) 的最大值为\fillinblank{}.

\bitmapfigure[max width=0.42\linewidth]{img/2026/shanghai_spring/q10_fig1.png}
\end{problem}

\begin{answer}
\(-\dfrac{39}{32}\)
\end{answer}

\begin{solution}
由 \(BD=DE=EC\)，设 \(\overrightarrow{AD}=\bs{u}\)，\(\overrightarrow{AE}=\bs{v}\). 可写出
\(\overrightarrow{AB}=2\bs{u}-\bs{v}\)，\(\overrightarrow{AC}=-\bs{u}+2\bs{v}\).
于是
\(\overrightarrow{AB}\cdot\overrightarrow{AC}=-2|\bs{u}|^2+5\bs{u}\cdot\bs{v}-2|\bs{v}|^2\).
由 \(AD=1\) 且 \(\angle DAE=\pi/3\)，令 \(|AE|=t>0\)，则
\(\overrightarrow{AB}\cdot\overrightarrow{AC}=-2t^2+\frac52t-2=-2\left(t-\frac58\right)^2-\frac{39}{32}\).
故最大值为 \(-\dfrac{39}{32}\).
\end{solution}

\begin{problem}
已知椭圆 \(\Gamma_1:\dfrac{x^2}{a^2}+y^2=1\)（\(a>1\)） 与椭圆 \(\Gamma_2:\dfrac{y^2}{b^2+2}+\dfrac{x^2}{b^2}=1\) 相交于 \(A\)，\(B\)，\(C\)，\(D\) 四点，且 \(\Gamma_1\) 和 \(\Gamma_2\) 的四个焦点在同一个圆上，则 \(b^2=\fillinblank{}\).
\end{problem}

\begin{answer}
\(3\)
\end{answer}

\begin{solution}
由两椭圆的对称性，四个焦点所在圆的圆心为原点. 两椭圆的半焦距分别为 \(\sqrt{a^2-1}\) 和 \(\sqrt{(b^2+2)-b^2}=\sqrt2\)，故 \(a^2-1=2\)，即 \(a^2=3\). 四个交点也在该圆上，故满足 \(x^2+y^2=2\). 代入 \(\Gamma_2\) 得
\(\frac{y^2}{b^2+2}+\frac{x^2}{b^2}=1\)，\(x^2+y^2=2\)
从而 \(b^2=3\).
\end{solution}

\begin{problem}
有一个油壶，壶身视为圆柱，壶嘴视为直线且不计容积，壶底直径 \(6.4\text{ 厘米}\)，壶身高 \(16\text{ 厘米}\)，壶内油液面高 \(12.1\text{ 厘米}\)，壶嘴长 \(13.4\text{ 厘米}\)，与壶身夹角 \(30^\circ\)，壶嘴最低点距壶底 \(3\text{ 厘米}\). 将壶身向壶嘴方向至少转\fillinblank{}\(\text{ 度}\)可使油倒出（精确到 \(0.01^\circ\)）.

\bitmapfigure[max width=0.55\linewidth]{img/2026/shanghai_spring/q12_fig1.png}
\end{problem}

\begin{answer}
\(14.2^\circ\)
\end{answer}

\begin{solution}
\solutionfigure{\bitmapfigure[max width=0.55\linewidth]{img/2026/shanghai_spring/q12_solution_fig1.png}}

设壶嘴最低点，最高点分别是 \(E\)，\(F\). 图中圆柱轴截面为矩形，距离点 \(E\) 最近的顶点是 \(A\)，另外三个顶点分别为 \(B\)，\(C\)，\(D\). 当水平液面经过点 \(F\) 时，可将油倒出.

设倾斜角为 \(\theta\). 当液面经过点 \(D\) 时，\(\theta=\arctan\dfrac{3.9}{6.4}\approx31.36^\circ\). 先考虑液面不超过点 \(D\)，即 \(\theta\in\left[0,\arctan\dfrac{3.9}{6.4}\right]\) 的情况. 设液面与 \(AD\)，\(CB\) 分别交于点 \(G\)，\(H\)，设 \(GH\) 的中点为 \(M\)，过 \(M\) 作 \(AD\) 的垂线，垂足为 \(N\). 则 \(MN=\dfrac12AB=3.2\)，\(\angle GMN=\theta\)，所以 \(GN=MN\tan\theta=3.2\tan\theta\). 因为 \(AN=12.1\)，所以 \(GE=GN+AN-AE=3.2\tan\theta+9.1\).

因为 \(\angle EGF=\angle NMG+\angle MNG=90^\circ+\theta\)，所以 \(\angle F=180^\circ-\angle FEG-\angle EGF=60^\circ-\theta\). 在 \(\triangle EGF\) 中，由正弦定理，
\[
\frac{3.2\tan\theta+9.1}{\sin(60^\circ+\theta)}
=\frac{13.4}{\sin(90^\circ+\theta)}
=\frac{13.4}{\cos\theta}
\]
即
\[
3.2\sin\theta+9.1\cos\theta
=13.4\sin(60^\circ+\theta)
=6.7\sqrt3\cos\theta-6.7\sin\theta
\]
所以 \(99\sin\theta=(67\sqrt3-91)\cos\theta\)，即
\[
\theta=\arctan\frac{67\sqrt3-91}{99}\approx14.2^\circ
\]
因为 \(14.2^\circ<31.36^\circ\)，所以至少将壶身倾斜 \(14.2^\circ\) 即可将油倒出.
\end{solution}

\section{单选题}

\begin{problem}
下列数列中是等差数列也是等比数列的是（\quad）.

\choices
    {1，\(-1\)，1，\(-1\)，1}
    {1，2，3，4，5}
    {5，5，5，5，5}
    {1，2，3，5，7}
\end{problem}

\begin{answer}
C
\end{answer}

\begin{solution}
常数列既是等差数列，也是等比数列，故选 C.
\end{solution}

\begin{problem}
已知 \(x>y>1\)，则下列不等式恒成立的是（\quad）.

\choices
    {\(x>y^2\)}
    {\(xy>x+y\)}
    {\(x^2>y\)}
    {\(x+y>xy\)}
\end{problem}

\begin{answer}
C
\end{answer}

\begin{solution}
由 \(x>y>1\) 得 \(x^2>x>y\)，故 \(x^2>y\) 恒成立. 其余选项均可举反例否定.
\end{solution}

\begin{problem}
平移对称法在几何学中具有重要的应用. 设平面直角坐标系 \(xOy\) 中有一图形 \(\Omega\)，过 \(\Omega\) 内任意一点 \(P\) 作垂直于 \(x\) 轴的直线 \(l_P\)，满足 \(l_P\cap\Omega\) 为一线段. 现沿 \(l_P\) 方向平移这些线段，使得它们的中点均在 \(x\) 轴上，这样叫做平移对称法. 对于由 \(-x^2+x+1\)，\(-x^2-x\)，直线 \(x=0\) 和直线 \(x=1\) 围成的封闭图形 \(\Omega\)，对它进行一次平移对称，得到的图像大致为（\quad）.

\choices
    {\choicebitmap[width=0.15\paperwidth]{img/2026/shanghai_spring/q15_choice_a.png}}
    {\choicebitmap[width=0.15\paperwidth]{img/2026/shanghai_spring/q15_choice_b.png}}
    {\choicebitmap[width=0.15\paperwidth]{img/2026/shanghai_spring/q15_choice_c.png}}
    {\choicebitmap[width=0.15\paperwidth]{img/2026/shanghai_spring/q15_choice_d.png}}
\end{problem}

\begin{answer}
A
\end{answer}

\begin{solution}
原图形上下两条边分别是两条开口向下的抛物线，平移对称后应保持关于 \(x\) 轴的对应关系. 对照四个选项，只有 A 与“中点在 \(x\) 轴上”的平移对称结果一致，故选 A.
\end{solution}

\begin{problem}
对于函数 \(y=f(x)\)，\(x\in D\)，设 \(A_f=\{(x,y)\mid y\ge f(x),x\in D\}\). 对于点集 \(M\)，若存在 \((x_0,y_0)\in M\)，使得任取 \((x,y)\in M\)，总有 \(y\ge y_0\)，则称 \((x_0,y_0)\) 为“最低点”. 对于函数 \(y=f(x)\) 和 \(y=g(x)\)，以下说法中正确的是（\quad）.

\choices
    {若 \(y=f(x)\) 和 \(y=g(x)\) 都有最小值，则 \(A_f\cap A_g\) 有最低点}
    {若 \(A_f\cap A_g\) 有最低点，则 \(y=f(x)\) 和 \(y=g(x)\) 都有最小值}
    {若 \(y=f(x)\) 或 \(y=g(x)\) 有最小值，则 \(A_f\cup A_g\) 有最低点}
    {若 \(A_f\cup A_g\) 有最低点，则 \(y=f(x)\) 或 \(y=g(x)\) 有最小值}
\end{problem}

\begin{answer}
D
\end{answer}

\begin{solution}
若 \(A_f\cup A_g\) 有最低点，则该最低点必对应某个函数图像上的最低值点，因此 \(y=f(x)\) 或 \(y=g(x)\) 至少有一个有最小值. 其余三项都可构造反例否定，故选 D.
\end{solution}

\section{解答题}

\begin{problem}
某兴趣班共\(150\text{ 人}\)，年龄分布及兴趣爱好统计如下：

\begin{center}
\renewcommand{\arraystretch}{1.25}
\begin{tabular}{|c|c|c|c|c|}
\hline
年龄 & 剪纸 & 摄影 & 画画 & 人数 \\
\hline
\([25,35)\) & & 8 & & 45 \\
\hline
\([35,45)\) & & 10 & & 55 \\
\hline
\([45,55)\) & & 6 & & 50 \\
\hline
\end{tabular}
\end{center}

\begin{enumerate}
\item 现采用分层抽样抽取\(30\text{ 人}\)，其中抽到年龄在\([25,35)\text{ 岁}\)的有多少人？
\item 该兴趣班\(150\text{ 人}\)的平均年龄是多少？
\item 现从\(150\text{ 人}\)中任意抽选\(1\text{ 人}\)，记抽到的学员年龄在\([35,45)\)为事件\(A\)，记抽到学员爱好摄影为事件\(B\). 事件\(A\)与\(B\)是否独立？请说明理由.
\end{enumerate}
\end{problem}

\begin{solution}
\begin{enumerate}
\item \([25,35)\) 岁的人数占 \(45/150\)，故抽到的人数为 \(30\times \frac{45}{150}=9\).

\item 将各组年龄视为组中值 \(30\)，\(40\)，\(50\)，则平均年龄为
\(\bar x=\frac{30\times45+40\times55+50\times50}{150}=\frac{121}{3}\).

\item 由表得
\(P(A)=\frac{55}{150}=\frac{11}{30}\)，\(P(B)=\frac{8+10+6}{150}=\frac{4}{25}\)，\(P(AB)=\frac{10}{150}=\frac{1}{15}\). 且 \(P(A)P(B)=\frac{11}{30}\cdot\frac{4}{25}\ne \frac{1}{15}\)，
所以 \(A\) 与 \(B\) 不独立.
\end{enumerate}
\end{solution}

\begin{problem}
如图所示正四棱台 \(ABCD-A_1B_1C_1D_1\)，其中 \(AB=4\)，\(A_1B_1=2\).
\begin{enumerate}
\item 当 \(AA_1=2\) 时，求 \(AA_1\) 和平面 \(A_1B_1C_1D_1\) 所成角；
\item 证明：\(AA_1\parallel\) 平面 \(BC_1D\)；若棱台高为 3，求三棱锥 \(A_1-BC_1D\) 的体积.
\end{enumerate}

\bitmapfigure[max width=0.42\linewidth]{img/2026/shanghai_spring/q18_fig1.png}
\end{problem}

\begin{solution}
\solutionfigure{\bitmapfigure[max width=0.55\linewidth]{img/2026/shanghai_spring/q18_solution_fig1.png}}

\begin{enumerate}
\item 过 \(A_1\) 作 \(A_1H\perp\) 平面 \(ABCD\) 于 \(H\)，连接 \(AH\). 过 \(H\) 分别作 \(HE\perp AB\) 于 \(E\)，\(HF\perp AD\) 于 \(F\)，连接 \(A_1E\)，\(A_1F\). 因为 \(AB\subset\) 平面 \(ABCD\)，所以 \(A_1H\perp AB\). 又 \(A_1H\cap HE=H\)，且 \(A_1H\)，\(HE\subset\) 平面 \(A_1HE\)，故 \(AB\perp\) 平面 \(A_1HE\)，从而 \(A_1E\perp AB\). 因此 \(AE=\dfrac{4-2}{2}=1\)，同理 \(AF=1\)，四边形 \(AEHF\) 为正方形，所以 \(AH=\sqrt2\). \(AH\) 是 \(A_1A\) 在平面 \(ABCD\) 上的投影，又平面 \(ABCD\parallel\) 平面 \(A_1B_1C_1D_1\)，故所成角为 \(\angle A_1AH\)，且
\[
\cos\angle A_1AH=\frac{AH}{AA_1}=\frac{\sqrt2}{2}
\]
所以所成角为 \(45^\circ\).

\item 连接 \(AC\)，\(BD\) 交于 \(O\)，连接 \(A_1C_1\)，\(B_1D_1\) 交于 \(O_1\). 上下底面为正方形，由正棱台性质，\(A_1C_1\parallel AC\)，且 \(A_1C_1=2\sqrt2\)，\(AC=4\sqrt2\)，\(AO=\dfrac12AC=A_1C_1\). 因此四边形 \(A_1C_1OA\) 为平行四边形，\(AA_1\parallel OC_1\). 又 \(AA_1\not\subset\) 平面 \(BC_1D\)，\(OC_1\subset\) 平面 \(BC_1D\)，所以 \(AA_1\parallel\) 平面 \(BC_1D\).

由正棱台性质，\(OO_1\) 与上下底面均垂直，故 \(OO_1=3\). 又 \(OO_1\perp BD\)，\(AC\perp BD\)，且 \(AC\cap OO_1=O\)，所以 \(BD\perp\) 平面 \(A_1OC_1\). 因此三棱锥体积可拆分为两个小三棱锥的体积之和：
\[
V_{A_1-BC_1D}=V_{B-A_1OC_1}+V_{D-A_1OC_1}
=\frac13OB\cdot S_{\triangle A_1OC_1}+\frac13OD\cdot S_{\triangle A_1OC_1}
=\frac13BD\cdot S_{\triangle A_1OC_1}=8
\]
\end{enumerate}

\end{solution}

\begin{problem}
已知 \(f(x)=\sin(\omega x+\varphi)\)（\(\omega>0\)，\(0<\varphi<\pi\)）.

\begin{enumerate}
\item 当 \(\omega=2\)，\(f\!\left(\dfrac{\pi}{12}\right)=0\)，求函数 \(f(x)\) 在 \(x=0\) 处的切线方程；
\item 若函数 \(f(x)\) 的最小正周期为 \(3\pi\)，且 \(f(x)=\dfrac{\sqrt2}{2}\) 在区间 \([0,2026\pi)\) 上恰好有 1351 个解，求 \(\varphi\) 的取值范围.
\end{enumerate}
\end{problem}

\begin{solution}
\begin{enumerate}
\item 当 \(\omega=2\) 时，\(f(x)=\sin(2x+\varphi)\). 由 \(f\left(\dfrac{\pi}{12}\right)=0\) 得 \(\sin\left(\dfrac{\pi}{6}+\varphi\right)=0\)，所以 \(\dfrac{\pi}{6}+\varphi=k\pi\). 结合 \(0<\varphi<\pi\)，得 \(\varphi=\dfrac{5\pi}{6}\)，即 \(f(x)=\sin\left(2x+\dfrac{5\pi}{6}\right)\). 因而 \(f'(0)=2\cos\dfrac{5\pi}{6}=-\sqrt3\)，且 \(f(0)=\dfrac12\)，切线方程为 \(y=-\sqrt3x+\dfrac12\).

\item 最小正周期为 \(3\pi\)，故 \(\omega=\dfrac23\)，于是 \(f(x)=\sin\left(\dfrac23x+\varphi\right)\). 令 \(t=\dfrac23x+\varphi\)，当 \(x\in[0,2026\pi)\) 时，\(t\in\left[\varphi,\dfrac{4052\pi}{3}+\varphi\right)\). 由 \(\sin t=\dfrac{\sqrt2}{2}\)，得 \(t=\dfrac{\pi}{4}+2k\pi\) 或 \(t=\dfrac{3\pi}{4}+2k\pi\).

当 \(0<\varphi\le\dfrac{\pi}{4}\) 时，要有 \(1351\) 个解，需满足
\(\frac{\pi}{4}+2\pi\cdot675<\frac{4052\pi}{3}+\varphi\)，\(\frac{3\pi}{4}+2\pi\cdot675\ge\frac{4052\pi}{3}+\varphi\)
解得 \(0<\varphi\le\dfrac{\pi}{12}\). 当 \(\dfrac{\pi}{4}<\varphi\le\dfrac{3\pi}{4}\) 时，同理需满足
\(\frac{3\pi}{4}+2\pi\cdot675<\frac{4052\pi}{3}+\varphi\)，\(\frac{9\pi}{4}+2\pi\cdot675\ge\frac{4052\pi}{3}+\varphi\)
解得 \(\dfrac{\pi}{4}<\varphi\le\dfrac{3\pi}{4}\). 当 \(\dfrac{3\pi}{4}<\varphi<\pi\) 时无解. 综上，\(0<\varphi\le\dfrac{\pi}{12}\) 或 \(\dfrac{\pi}{4}<\varphi\le\dfrac{3\pi}{4}\).
\end{enumerate}
\end{solution}

\begin{problem}
已知双曲线 \(\Gamma:\dfrac{x^2}{2}-\dfrac{y^2}{2}=1\)，过点 \(M(m,0)\) 作不垂直于 \(x\) 轴的直线 \(l\) 交双曲线 \(\Gamma\) 于 \(A\)，\(B\) 两点.

\begin{enumerate}
\item 求双曲线 \(\Gamma\) 的离心率；
\item 若点 \(A(\sqrt3,1)\)，点 \(B\) 在双曲线的右支上，且 \(B\) 是 \(AM\) 的中点，求直线 \(l\) 的斜率；
\item 若 \(m>0\)，\(F_1\)，\(F_2\) 分别是双曲线 \(\Gamma\) 的左右焦点，\(A'\) 是 A 关于 \(y\) 轴的对称点，若存在直线 \(l\) 使得 \(\overrightarrow{F_1A'}\cdot\overrightarrow{F_2B}=0\)，求 \(m\) 的取值范围.
\end{enumerate}
\end{problem}

\begin{solution}
\begin{enumerate}
\item 由 \(\dfrac{x^2}{2}-\dfrac{y^2}{2}=1\) 可知 \(a^2=2\)，\(b^2=2\)，故 \(c^2=a^2+b^2=4\)，离心率 \(e=\frac ca=\sqrt2\).

\item 设 \(B\) 为 \(AM\) 的中点，则
\(B\left(\frac{\sqrt3+m}{2},\frac12\right)\).
又 \(B\in\Gamma\)，故
\(\frac{(\sqrt3+m)^2}{8}-\frac{1}{8}=1\)，
得 \(\sqrt3+m=3\)，即 \(m=3-\sqrt3\).
因此直线 \(l\) 的斜率为
\(k=\frac{1-0}{\sqrt3-(3-\sqrt3)}=\frac{2\sqrt3+3}{3}\).

\item 当直线斜率为 \(0\) 时，\(\overrightarrow{F_1A'}\) 与 \(\overrightarrow{F_2B}\) 共线，不符合题意. 当斜率不为 \(0\) 时，设 \(l:x=ny+m\)，\(n\ne0\). 设 \(A(x_1,y_1)\)，\(B(x_2,y_2)\)，则 \(A'(-x_1,y_1)\).
联立双曲线与直线得
\[
(n^2-1)y^2+2nmy+m^2-2=0
\]
故 \(\Delta=4(m^2+2n^2-2)>0\)，且 \(n^2-1\ne0\)，并有
\(y_1+y_2=\frac{-2nm}{n^2-1}\)，\(y_1y_2=\frac{m^2-2}{n^2-1}\)
由 \(F_1(-2,0)\)，\(F_2(2,0)\) 及 \(\overrightarrow{F_1A'}\cdot\overrightarrow{F_2B}=0\)，化简得 \(n^2=m^2-2m+1\). 代入 \(\Delta>0\) 得 \(3m^2-4m>0\)，结合 \(m>0\) 得 \(m>\dfrac43\). 又 \(n^2-1\ne0\) 给出 \(m\ne2\). 因此 \(m\in\left(\dfrac43,2\right)\cup(2,+\infty)\).
\end{enumerate}
\solutionfigure{\bitmapfigure[max width=0.55\linewidth]{img/2026/shanghai_spring/q20_solution_fig1.png}}
\end{solution}

\begin{problem}
设 \(f(x)\) 是定义在 \(\R\) 上的函数. 定义性质 \(P\)：若对任意 \(x_1\)，\(x_2\in\R\)，当 \(|x_1|<|x_2|\) 时，\(f(x_1)<f(x_2)\)，则称函数 \(f(x)\) 具有“性质 \(P\)”.

\begin{enumerate}
\item 判断函数 \(f(x)=e^x\) 是否具有“性质 \(P\)”；
\item 若分段函数
\(
f(x)=
\begin{cases}
ax,&x\le0,\\
x+b,&x>0
\end{cases}
\)
具有“性质 \(P\)”，求所有满足条件的实数 \(a\)，\(b\)；
\item 已知 \(f(x)\) 的值域为 \([0,1)\)，且在 \([0,+\infty)\) 上是严格增函数，证明：\(f(x)\) 是偶函数的充要条件是 \(f(x)\) 具有“性质 \(P\)”.
\end{enumerate}
\end{problem}

\begin{solution}
\begin{enumerate}
\item 函数 \(f(x)=e^x\) 不具有性质 \(P\). 取 \(x_1=-2\)，\(x_2=-3\)，则 \(|x_1|<|x_2|\)，但 \(e^{x_1}=e^{-2}>e^{-3}=e^{x_2}\)，与性质 \(P\) 矛盾.

\item 因为函数具有性质 \(P\)，取 \(x_1=0\)，有 \(f(0)<f(x_2)\)（\(x_2\ne0\)）. 当 \(x_2>0\) 时，\(0<x_2+b\)，故 \(b\ge0\)；当 \(x_2<0\) 时，\(0<ax_2\)，故 \(a<0\). 若 \(b>0\)，取 \(x_2\in\left[\dfrac{b}{a},0\right)\)，则 \(f(x_2)\le f\left(\dfrac{b}{a}\right)=b\). 取 \(x_1=-\dfrac{x_2}{2}>0\)，则 \(|x_1|<|x_2|\)，但 \(f(x_1)=x_1+b>b\ge f(x_2)\)，矛盾，故 \(b=0\).

此时
\[
f(x)=
\begin{cases}
ax,&x\le0,\\
x,&x>0.
\end{cases}
\]
若 \(x_1>0\)，\(x_2<0\)，\(x_1<-x_2\)，则 \(\dfrac{x_1}{x_2}\ge-1\). 由 \(f(x_1)<f(x_2)\) 得 \(x_1<ax_2\)，即 \(\dfrac{x_1}{x_2}>a\)，故 \(a\le-1\). 若 \(x_1<0\)，\(x_2>0\)，\(-x_1<x_2\)，则 \(\dfrac{x_2}{x_1}\le-1\). 由 \(f(x_1)<f(x_2)\) 得 \(ax_1<x_2\)，即 \(\dfrac{x_2}{x_1}<a\)，故 \(a\ge-1\). 因此 \(a=-1\)，\(b=0\).

\item 充分性：若 \(f(x)\) 具有性质 \(P\)，则对任意 \(x_0>0\)，若 \(f(x_0)\ne f(-x_0)\)，不妨设 \(f(x_0)>f(-x_0)\). 记 \(M=f(x_0)\)，\(m=f(-x_0)\)，则 \(0\le m<M<1\). 当 \(|x|<x_0\) 时，性质 \(P\) 给出 \(f(x)<m\)；当 \(|x|>x_0\) 时，给出 \(f(x)>M\). 因而对任意 \(x\in\R\)，\(f(x)\ne\dfrac{M+m}{2}\)，与值域为 \([0,1)\) 矛盾. 故 \(f(x_0)=f(-x_0)\)，\(f\) 为偶函数.

必要性：若 \(f\) 是偶函数，则 \(f(x)=f(|x|)\). 当 \(|x_1|<|x_2|\) 时，由 \(f\) 在 \([0,+\infty)\) 上严格增，得 \(f(|x_1|)<f(|x_2|)\)，从而 \(f(x_1)<f(x_2)\)，即 \(f\) 具有性质 \(P\). 因此 \(f(x)\) 是偶函数的充要条件是 \(f(x)\) 具有性质 \(P\).
\end{enumerate}
\end{solution}
